\subsection{معرفی و انگیزه}
این مکانیزم بر مبنای مبانی نظری پیشرفته در ادبیات توجه و فشرده‌سازی \cite{vaswani2017attention} توسعه یافته است.

افزایش طول توالی در مدل‌های زبانی بزرگ، محاسبات مکانیزم توجه را به دلیل پیچیدگی زمانی و مکانی $\mathcal{O}(N^2)$ با چالش‌های حاد مواجه می‌سازد. ماژول نوآورانه \lr{Aegis-Spectron-FP8} با تلفیق سه پارادایم محاسباتی ناهمگن به همراه کوانتیزاسیون هشت‌بیتی با ممیز شناور \lr{(FP8)} و مسیریابی مبتنی بر آنتروپی طیفی طراحی شده است تا مصرف حافظه گرافیکی را به کمینه ممکن رسانده و گلوگاه ارتباطی گذرگاه حافظه را بشکند.

ایده محوری این ماژول بر پایه تفکیک فضای توجه به سه شاخه مستقل استوار است:
\begin{enumerate}
    \item \textbf{شاخه توجه متحرک محلی (\lr{Local Sliding-Window}):} با پیچیدگی $\mathcal{O}(N \cdot W)$ جهت پردازش دقیق کلمات همجوار.
    \item \textbf{شاخه خطی کِرنلی کم‌مرتبه (\lr{Global ELU+1 Linear Attention}):} با پیچیدگی $\mathcal{O}(N \cdot r \cdot d)$ برای ثبت الگوهای معنایی در پهنه متن.
    \item \textbf{شاخه بازگشتی خطی حالت‌فضا (\lr{Linear State Recurrence}):} با پیچیدگی خطی $\mathcal{O}(N \cdot D)$ بر پایه معادلات بازگشتی جهت انتقال تاریخچه متنی.
\end{enumerate}
ترکیب این شاخه‌ها از طریق مسیریاب آنتروپی طیفی (\lr{Spectron Spectral Entropy Router}) و اعمال تبدیل‌های فشرده‌سازی عددی \lr{FP8} در مراحل پیشرو صورت می‌پذیرد.
